\chapter{Phlebotomy Considerations in Children}

\section*{Overview}
Phlebotomy in children is the procedure of collecting blood samples from pediatric patients using age-appropriate techniques that minimize pain, distress, and blood volume loss.

\section*{Alternative Names}
Pediatric phlebotomy, children's blood draw, child venipuncture, pediatric blood collection, infant blood sampling


\section*{Principle}
Pediatric phlebotomy uses age-appropriate techniques: heel stick for neonates (warmed heel, lancet depth <2 mm), scalp vein or jugular vein for infants, and antecubital venipuncture for older children. Maximum blood draw volume is calculated as 1-3% of total blood volume (80 mL/kg) per collection.

\section*{History}
Pediatric phlebotomy techniques were developed in the mid-20th century as awareness grew that children require different approaches than adults. The use of topical anesthetics and smaller-gauge needles became standard.

\section*{Why It Is Done}
This section addresses the practical challenges of collecting blood from children of different ages. Understanding the best approach for each age group helps your doctor or phlebotomist obtain accurate results while minimizing pain and distress for you and your child. It answers the question: what is the safest and least stressful way to collect blood from a child of this age?

\section*{Is This a Primary or Secondary Test or Procedure}
Informational guideline.

\section*{How It Is Performed}

\begin{longtable}{@{}ll@{}}
\toprule
\textbf{Age Group} & \textbf{Preferred Collection Sites} \\
\midrule
Newborns (0--1 month) & Heel prick, jugular vein, femoral vein, PICC line \\
Infants (1--12 months) & Jugular vein, femoral vein, saphenous vein \\
Toddlers (1--3 years) & Antecubital fossa (with distraction techniques) \\
Children (3+ years) & Antecubital fossa preferred \\
\bottomrule
\end{longtable}

Key principles for pediatric blood collection:
\begin{itemize}
\item Minimize the volume of blood drawn --- use pediatric tubes
\item Use distraction techniques (toys, videos, storytelling) for young children
\item Consider point-of-care testing when available to reduce the need for venous draws
\item Use non-invasive alternatives when possible (transcutaneous bilirubin instead of blood bilirubin, pulse oximetry instead of arterial blood gas)
\item Comfort and reassure your child throughout the process
\end{itemize}

\section*{Preparation}
Psychological preparation of child and parents; application of topical anesthetic (EMLA cream) 1 hour prior.

\section*{Contraindications and Precautions}
Avoid crying or hyperventilation before collection to prevent transient shifts in blood gas parameters.

\section*{Risks}
Bruising, minor bleeding, or psychological distress/needle phobia.

\section*{Aftercare and Recovery}
Hold pressure for 3-5 minutes; reward the child for cooperation; monitor for hematoma.

\section*{Outcome and Follow-up}
Results should be interpreted using age-specific pediatric reference ranges.

\section*{Expected Results}
Successful blood collection with minimal hemolysis or trauma.

\section*{Follow-up}
Monitoring serial labs if necessary.


\section*{References}
\begin{enumerate}
  \item MedlinePlus. Pediatric phlebotomy. U.S. National Library of Medicine; 2025.
  \item Current Medical Diagnosis and Treatment. McGraw-Hill; 2025.
\end{enumerate}

\clearpage
